\chapter{1977年江苏省农垦系统卷}

\begin{problem}
计算：

（1）\( \left[ 1\frac{1}{24} - \left( \frac{3}{8} + \frac{1}{6} - \frac{3}{4} \right) \times 24 \right] \div \left( -5 \right)^{0} \)；

（2）\( \lg \frac{5}{2} + \lg \frac{2}{5} - \lg 10^{3} \)。

\end{problem}

\begin{problem}
解方程 \( \frac{1}{x + 2} + \frac{4x}{x^{2} - 4} - \frac{2}{x - 2} = 1 \)。

\end{problem}

\begin{problem}
用长 \( 60 \) 米的篱笆围成一矩形养鸡场地，一面利用原有旧墙，问矩形的长和宽各多少米才能使所围成的面积最大？

\end{problem}

\begin{problem}
把极坐标方程 \( \rho = 2 \cos \theta \) 化成直角坐标方程，并说出它是什么曲线？

\end{problem}

\begin{problem}
在锐角 \( \triangle ABC \) 中，已知 \( \sin A = \frac{3}{5} \)，\( \cos B = \frac{5}{13} \)，求 \( \cos C \) 的值。

\end{problem}

\begin{problem}
三个数成等差数列，且它们的和是 \( 9 \)，如果把这个数列依次加上 \( 1 \)、\( 1 \)、\( 3 \)，那么所得的数又成等比数列。求这三个数。

\end{problem}

\begin{problem}
求以椭圆 \( \frac{x^{2}}{49} + \frac{y^{2}}{25} = 1 \) 的焦点和长轴的端点作为顶点和焦点的双曲线方程。

\end{problem}

\begin{problem}
等腰三角形 \( ABC \) 的顶角 \( A \) 为 \( 100^{\circ} \)，底角 \( B \) 的平分线交腰 \( AC \) 于 \( D \)，求证：\( BC = AD + BD \)。
\end{problem}

